% =====================================================================
%  RESUMEN DE LA CARACTERIZACIÓN DE LAS TECNOLOGÍAS
%  PROXY-CACHÉ Y ALMACENAMIENTO DE ARCHIVOS
%
%  Marco operacional para una solución basada en proxy-caché
%  en la infraestructura del Grupo de Investigación GRID
%  Universidad del Quindío — 2025
%
%  Compilar:  pdflatex main.tex   (dos veces, por las referencias)
%         o:  latexmk -pdf main.tex
%
%  Para agregar una sección nueva: crea el archivo en secciones/
%  y añade su \input al final de la lista de abajo.
% =====================================================================
\documentclass[11pt,a4paper]{article}

\usepackage{mathptmx}                 % tipografía tipo Times (opcional)
% =====================================================================
%  PREÁMBULO — Análisis DAR
%  Configuración de paquetes, colores, tablas y leyendas.
% =====================================================================

% ---------- Codificación e idioma -----------------------------------
\usepackage[utf8]{inputenc}
\usepackage[T1]{fontenc}
% Si tu distribución tiene el idioma instalado (Overleaf/TeX Live full),
% descomenta la línea siguiente para partición silábica en español:
% \usepackage[spanish,es-nodecimaldot]{babel}

% ---------- Geometría de página --------------------------------------
\usepackage[a4paper,margin=2.5cm]{geometry}

% ---------- Tablas ----------------------------------------------------
\usepackage{array}
\usepackage{makecell}
\usepackage[table]{xcolor}
\usepackage{ragged2e}
\usepackage{float}        % habilita el especificador [H] (tabla inquebrable)
\usepackage{graphicx}
\usepackage{caption}
\usepackage{fancyhdr}
\setlength{\headheight}{13.6pt}

% ---------- Enlaces ---------------------------------------------------
\usepackage{hyperref}
\hypersetup{
  colorlinks=true,
  linkcolor=black,
  citecolor=black,
  urlcolor=blue,
  pdftitle={Análisis DAR — Proxy-Caché y Almacenamiento de Archivos},
  pdfauthor={Santiago Quintero Uribe, Juan Manuel Perdomo Cárdenas, Andrés Felipe Zúñiga Zuluaga}
}

% ---------- Numeración de página estilo APA 7 ------------------------
% APA 7 para trabajos estudiantiles usa solo el número de página en el
% encabezado superior derecho. No se incluye running head salvo que se exija.
\pagestyle{fancy}
\fancyhf{}
\fancyhead[R]{\thepage}
\renewcommand{\headrulewidth}{0pt}

% ---------- Colores institucionales / de tabla ------------------------
\definecolor{gridhead}{HTML}{C5E0B3}   % verde claro de encabezado
\definecolor{gridrow}{HTML}{CCCCCC}    % gris de filas alternas
\definecolor{verdeuq}{HTML}{009B3A}    % verde Universidad del Quindío

% =====================================================================
%  LEYENDAS (captions) DE TABLAS
% =====================================================================
% position=top     -> el paquete caption calcula el espaciado asumiendo
%                     que el \caption va ENCIMA del contenido.
% skip=6pt         -> separación entre la leyenda y la tabla.
% justification    -> leyenda centrada si cabe en una línea, justificada
%                     si ocupa varias.
\captionsetup[table]{
  name=Tabla,
  labelfont=bf,
  labelsep=period,
  position=top,
  skip=6pt,
  font=small,
  justification=centering,
  singlelinecheck=true
}

% =====================================================================
%  APARIENCIA DE LAS TABLAS
% =====================================================================
\renewcommand{\arraystretch}{1.25}
\setlength{\tabcolsep}{4pt}

% Espacio antes/después de cada entorno table colocado con [H]
\setlength{\intextsep}{12pt plus 3pt minus 3pt}

% Tolerancia de página: evita que LaTeX estire en exceso el texto
% cuando una tabla se empuja a la página siguiente.
\raggedbottom


\begin{document}

% ---------- Portada --------------------------------------------------
% =====================================================================
%  PORTADA — Análisis DAR
%  Edita únicamente los datos del bloque siguiente.
% =====================================================================

\newcommand{\tipoDocumento}{Análisis DAR (Decision Analysis and Resolution)
de las Tecnologías Proxy-Caché y Almacenamiento de Archivos}

\newcommand{\tituloProyecto}{MARCO OPERACIONAL PARA UNA SOLUCIÓN BASADA EN
PROXY-CACHÉ EN LA INFRAESTRUCTURA DEL GRUPO DE INVESTIGACIÓN GRID DE LA
UNIVERSIDAD DEL QUINDÍO}

\newcommand{\autorUno}{Santiago Quintero Uribe}
\newcommand{\autorDos}{Juan Manuel Perdomo Cárdenas}
\newcommand{\autorTres}{Andrés Felipe Zúñiga Zuluaga}

\newcommand{\universidad}{Universidad del Quindío}
\newcommand{\facultad}{Facultad de Ingeniería}
\newcommand{\programa}{Ingeniería de Sistemas y Computación}
\newcommand{\ciudadAnio}{Armenia, Quindío, 2025}
\newcommand{\rutaLogo}{logos/logo_uniquindio.png}

% ---------------------------------------------------------------------
\begin{titlepage}
\centering
\bfseries

\vspace*{1.5cm}

{\large \tipoDocumento \par}

\vspace{1.4cm}

{\large \tituloProyecto . \par}

\vspace{1.8cm}

{\normalsize \autorUno \par}
\vspace{0.6em}
{\normalsize \autorDos \par}
\vspace{0.6em}
{\normalsize \autorTres \par}

\vfill

\includegraphics[height=4.2cm]{\rutaLogo}

\vfill

{\normalsize
\universidad \par
\facultad \par
\programa \par
\ciudadAnio \par}

\vspace*{0.8cm}
\end{titlepage}


% ---------- Tabla de contenido ---------------------------------------
\renewcommand{\contentsname}{Tabla de contenido}
\tableofcontents
\clearpage

% Si quieres además un índice de tablas, descomenta estas dos líneas:
% \listoftables
% \clearpage

% ---------- Contenido ------------------------------------------------
%% ====================================================================
%%  Metodología de evaluación
%% ====================================================================
\section{Metodología de evaluación}\label{sec:metodologia}

La metodología de evaluación aplicada para la elección de las
tecnologías de proxy-caché y almacenamiento de archivos de gran tamaño
fue DAR (Decision Analysis and Resolution) del modelo CMMI (Capability
Maturity Model Integration) (CMMI Institute, 2010). Este referente
metodológico permitió evaluar las necesidades, problemas y oportunidades
(NPO) del grupo de investigación GRID a través de un proceso
estructurado que consideró múltiples alternativas, criterios de
evaluación bien definidos y un análisis comparativo cuantitativo. En
este caso, se identificaron y analizaron 14 tecnologías proxy-caché
como: HAProxy, NGINX, Envoy, Varnish, Traefik, Caddy, Apache Traffic
Server, Kong, Squid, Istio, Linkerd, K8Sidecar, Proxy-Service y nuster y
22 tecnologías de almacenamiento incluyendo Ceph, Lustre, JuiceFS, DAOS,
SeaweedFS, HDFS, MinIO, Longhorn, entre otras; aplicando criterios
ponderados como rendimiento técnico (25 \%), escalabilidad (20 \%),
soporte y mantenimiento (20 \%), adopción y madurez (10 \%),
observabilidad (10 \%), adaptación al GRID y riesgo tecnológico (5 \%).
Mediante el uso de matrices de evaluación NPO y puntuación DAR, se logró
visualizar las fortalezas y debilidades de cada opción, facilitando una
selección alineada con los objetivos estratégicos del GRID. Así, el uso
de DAR no solo aportó transparencia al proceso, sino también
trazabilidad y justificación técnica frente a una decisión crítica para
la arquitectura de infraestructura del grupo de investigación.

%% ====================================================================
%%  Criterios de evaluación
%% ====================================================================
\section{Criterios de evaluación}\label{sec:criterios}

\subsection{Necesidades}

\subsubsection{Repositorio local archivos grandes}

Esta métrica evalúa la capacidad de la tecnología de almacenamiento para gestionar objetos de gran tamaño, como imágenes ISO y máquinas virtuales, de forma nativa y persistente. Una calificación de 3 requiere soporte nativo para objetos superiores a 1 GB con hit ratio mayor al 80 \%, replicación o erasure coding integrados, y persistencia garantizada documentada en benchmarks oficiales, lo cual evita descargas repetidas a través de la WAN y acelera los laboratorios del GRID.

\subsubsection{Priorización tráfico académico (QoS)}

Mide la capacidad de la herramienta para clasificar y priorizar el tráfico misional (docencia e investigación) frente al tráfico no misional. Una puntuación excelente exige mecanismos nativos avanzados de QoS con ACLs granulares, rate limiting por clase de tráfico, circuit breaking o priority routing, garantizando el ancho de banda necesario para las actividades académicas del grupo.

\subsubsection{Métricas uso red y proxy}

Evalúa si la solución ofrece métricas nativas de consumo de ancho de banda por usuario o dispositivo, y volumen por tipo de contenido. La calificación máxima requiere un endpoint /metrics compatible con Prometheus, stats API granular o CSI metrics nativo, permitiendo la toma de decisiones basada en datos (KPIs) sin depender de exporters externos obligatorios.

\subsubsection{Automatización certificados SSL}

Analiza el soporte nativo para la gestión automática de certificados firmados por entidades externas como Let's Encrypt. La puntuación más alta se otorga cuando la herramienta ofrece HTTPS/TLS automático ``by default'' con rotación integrada y sin necesidad de scripts externos ni intervención manual, eliminando advertencias de sitios inseguros y aumentando la confianza de la comunidad académica.

\subsection{Problemas}

\subsubsection{Tiempos de descarga deficientes}

Mide la latencia en la descarga de ISOs desde la red GRID hacia los usuarios finales. Una calificación de 3 exige sub-milisegundo en cache hit o tiempos documentados menores a 60 segundos para ISOs de 1 GB, con throughput masivo sostenido superior a 10 GB/s, resolviendo el cuello de botella de red actual de la infraestructura TI de la universidad.

\subsubsection{Uso no misional de red}

Evalúa la capacidad de identificar, filtrar y priorizar el tráfico de entretenimiento frente al académico. La puntuación máxima requiere ACLs granulares nativas, rate limiting por categoría y tagging o metadata nativo, permitiendo bloquear o priorizar tráfico no misional de forma declarativa o automática, evitando que obstaculice clases virtuales y laboratorios.

\subsubsection{Falta certificados SSL firmados}

Analiza si la herramienta permite eliminar las advertencias de sitios web inseguros del GRID mediante TLS nativo y compatibilidad confirmada con Let's Encrypt, incluyendo rotación automática de certificados sin intervención operativa, lo cual afecta directamente la confianza de la comunidad académica en los servicios internos.

\subsection{Oportunidades}

\subsubsection{Gobernanza del tráfico}

Mide la existencia de un framework declarativo y dinámico para el control de contenido que transita por la infraestructura. La calificación más alta requiere configuración declarativa nativa, versionado y runtime API dinámica (xDS, CRD Kubernetes, hot reload) que permita gobernanza en tiempo real sin downtime, otorgando control administrativo completo del tráfico.

\subsubsection{Optimización ancho de banda}

Evalúa los mecanismos para limitar el uso por usuario y priorizar el tráfico misional. Una puntuación de 3 exige rate limiting nativo por usuario o frontend, QoS por categoría y algoritmos de cola avanzados (prio, wrr, circuit breaking), maximizando el uso misional del ancho de banda y reduciendo costos WAN.

\subsubsection{Toma decisiones basada en datos}

Analiza la capacidad de generar KPIs declarados para monitoreo y reportes. La calificación máxima requiere dashboard nativo, métricas Prometheus completas y reportes automáticos integrados, permitiendo identificar patrones de uso y optimizar la infraestructura sin complementos críticos externos.

\subsection{Criterios Funcionales}

\subsubsection{Cache archivos grandes}

Evalúa la funcionalidad de caché HTTP optimizada para objetos mayores a 1 GB. Una calificación de 3 requiere caché nativa con hit ratio superior al 80 \% documentado, sub-milisegundo en cache hit o speedup confirmado de 100x+ en benchmarks oficiales, constituyendo el requisito funcional principal del proxy-caché para el GRID.

\subsubsection{QoS y priorización}

Mide los mecanismos para priorizar tráfico misional con granularidad por usuario, IP o contenido. La puntuación máxima exige QoS nativo múltiple con colas, rate limit, circuit breaking y priority routing, alineándose totalmente con los objetivos misionales del grupo de investigación.

\subsubsection{ACL listas blancas/negras}

Analiza el control de acceso por usuario, IP, contenido, header, método, payload, geolocalización, regex o JWT. Una calificación de 3 requiere ACLs nativas de granularidad máxima, permitiendo control total de acceso y fortaleciendo la seguridad perimetral del GRID.

\subsubsection{Proxy transparente/inverso/estándar}

Evalúa el soporte multimodal del proxy según la arquitectura requerida. La puntuación más alta exige soporte nativo de todos los modos (inverso, transparente TPROXY, estándar, passthrough, edge gateway) con documentación oficial de producción, garantizando flexibilidad arquitectónica.

\subsubsection{Gestión certificados digitales}

Mide el soporte para TLS, mTLS y rotación automática de certificados. Una calificación de 3 requiere TLS y mTLS nativos con rotación automática mediante SDS, auto-HTTPS o cert-manager nativo, sin downtime ni scripts externos, asegurando la seguridad en tránsito y la confianza de los usuarios.

\subsection{Criterios Técnicos}

\subsubsection{Escalabilidad almacenamiento}

Evalúa la capacidad de crecer horizontalmente sin reemplazar hardware. La puntuación máxima requiere auto-scaling horizontal nativo, replicación o erasure coding sin reemplazo de hardware y auto-healing (CRUSH, CSI, master/volume auto-scaling), asegurando el crecimiento futuro de los recursos del GRID.

\subsubsection{Volúmenes transferencia altos}

Mide la capacidad de soportar throughput masivo concurrente. Una calificación de 3 exige throughput sostenido superior a 10 GB/s documentado en benchmarks oficiales (ej. 200+ GB/s en HPC, $>$500K IOPS NVMe), garantizando el soporte a eventos masivos como talleres e instalación de laboratorios.

\subsubsection{Integración servicios existentes}

Analiza la compatibilidad con la infraestructura actual del GRID. La puntuación más alta requiere integración plug-and-play con infraestructura legacy y moderna, con más de 20 años en producción enterprise o académica, documentación extensa y market share probado, minimizando cambios arquitectónicos.

\subsubsection{Monitoreo tiempo real + alta disponibilidad}

Evalúa la capacidad de dashboard en tiempo real, alertas y alta disponibilidad. Una calificación de 3 requiere Prometheus nativo, health checks y HA multi-active (MDS/RGW, CSI, cluster nativo) con dashboard en tiempo real sin dependencias críticas externas, evitando pérdida de información y downtime.

\subsection{Criterios Monitoreo}

\subsubsection{Consumo ancho de banda por usuario}

Mide la granularidad de las métricas de bandwidth por usuario o dispositivo en tiempo real. La puntuación máxima exige métrica nativa granular con exporter oficial, CSI metrics o telemetry nativa, incluyendo histórico integrado, para garantizar el cumplimiento de políticas de uso y fairness en la red.

\subsubsection{Estadísticas tráfico por categoría}

Evalúa la clasificación automática del tráfico en académico, administrativo y no misional. Una calificación de 3 requiere clasificación automática nativa por categoría con tagging, ACL labels o analytics nativo por service/route, presentando top 10 categorías en dashboard para facilitar la priorización y gobernanza.

\subsubsection{Tiempos de respuesta promedio}

Mide la latencia promedio del proxy y del storage. La puntuación más alta exige sub-milisegundo en cache hit documentado en benchmarks oficiales con throughput sostenido (42k+ RPS), cumpliendo con el SLA de experiencia de usuario establecido para el GRID.

\subsubsection{Reportes}

Analiza la capacidad de generación de reportes periódicos automáticos. Una calificación de 3 requiere reportes automáticos nativos con exportación en múltiples formatos (PDF, CSV, email), distribución programada y analytics nativo, permitiendo la toma de decisiones administrativas sin depender de stacks externos como ELK o Loki.

\subsection{Criterios Gobernanza}

\subsubsection{Políticas formales}

Evalúa la existencia de un framework de políticas estructurado para los servicios del GRID. La puntuación máxima requiere framework declarativo nativo con versionado y runtime API dinámica (xDS, CRD Kubernetes, hot reload), permitiendo gobernanza estructurada sin downtime ni configuración ad-hoc.

\subsubsection{Diferenciar tráfico misional/no misional}

Mide la capacidad de clasificación automática y tagging del tráfico. Una calificación de 3 exige auto-classification nativa por metadata, ACL o tagging (header-based, labeling Kubernetes, RBAC filters) con aplicación automática de políticas diferenciadas, asegurando la alineación estratégica del GRID.

\subsubsection{Adaptación a cambios}

Analiza la capacidad de modificar políticas sin downtime. La puntuación más alta requiere hot reload, runtime API o xDS dinámico nativo, permitiendo modificar ACLs, backends, certificados y políticas en tiempo real, garantizando flexibilidad ante cambios de contexto operativo o académico.

%% ====================================================================
%%  Plantilla de evaluación
%% ====================================================================
\section{Plantilla de evaluación}\label{sec:plantilla}

A continuación, se presentan las plantillas de evaluación NPO
construidas para el análisis de las tecnologías proxy-caché y de
almacenamiento de archivos de gran tamaño. Las calificaciones siguen una
escala de 1 a 3, donde 1 es ``Muy limitado'', 2 es ``Aceptable'' y 3 es
``Excelente / Líder''.

En esta plantilla se presentan 6 indicadores que tienen un valor
ponderado con respecto a la puntuación final de cada una de las
herramientas; los indicadores que se tienen son evaluados con respecto a
las métricas que se habían definido en el formulario para la
caracterización de la organización. Cada una de estas métricas se
definen en la sección~\ref{sec:criterios}, criterios de evaluación, en donde se detalla la
importancia de cada una de estas en el contexto del proyecto.

A continuación, se le da un sentido a cada uno de los indicadores con
respecto a las métricas.

\begin{itemize}
\item \textbf{Rendimiento técnico (25\%):} Repositorio local de archivos grandes, priorización de tráfico académico, tiempos de descarga deficientes, optimización de ancho de banda, caché de archivos grandes, volúmenes de transferencia altos y tiempos de respuesta promedio.
\item \textbf{Escalabilidad (20\%):} Repositorio local de archivos grandes, priorización de tráfico académico (QoS), uso no misional de red, gobernanza del tráfico, QoS y priorización, ACL listas blancas/negras, escalabilidad de almacenamiento, políticas formales y diferenciar tráfico misional/no misional.
\item \textbf{Soporte y Mantenimiento (20\%):} Automatización de certificados SSL, falta de certificados SSL firmados, gestión de certificados digitales y adaptación a cambios.
\item \textbf{Adopción y Madurez (10\%):} Automatización de certificados SSL, proxy transparente/inverso/estándar e integración con servicios existentes.
\item \textbf{Observabilidad (10\%):} Uso no misional de red, métricas de uso de red y proxy, toma de decisiones basada en datos, monitoreo en tiempo real y alta disponibilidad, consumo de ancho de banda por usuario, estadísticas de tráfico por categoría y reportes.
\item \textbf{Adecuación al Contexto GRID (10\%):} Ajuste operativo de la herramienta en el contexto de operaciones del GRID y su escalabilidad.
\item \textbf{Riesgo Tecnológico (5\%):} Falta de certificados SSL firmados y sección para los estados proyectados para el año 2026 en la Matriz DAR.
\end{itemize}

\subsection{Cómo se calcula la puntuación de cada herramienta}\label{subsec:formalizacion}

Para que cualquier persona pueda reproducir o verificar los resultados,
en esta subsección se explica paso a paso cómo se pasa de las
calificaciones individuales de cada métrica a la puntuación DAR final de
una herramienta. El procedimiento consta de dos pasos: primero se calcula
el valor de cada indicador y después se combinan todos los indicadores en
una sola puntuación.

\subsubsection{Paso 1: obtener el valor de cada indicador}

Cada indicador agrupa varias métricas y todas ellas pesan lo mismo dentro
del indicador. Por eso, el valor del indicador es simplemente el
\textbf{promedio} de las calificaciones que la herramienta obtuvo en sus
métricas; es decir, se suman esas calificaciones y el resultado se divide
entre el número de métricas:

\[
  \text{Valor del indicador} \;=\;
  \frac{\text{suma de las calificaciones de sus métricas}}
       {\text{cantidad de métricas del indicador}}
\]

Como cada métrica se califica de 1 a 3, el promedio siempre queda también
entre 1 y 3. Esto es importante porque permite comparar indicadores entre
sí aunque unos agrupen más métricas que otros.

Según la lista anterior, cada indicador se calcula así:

\[
  \text{Rendimiento técnico} =
  \frac{\text{suma de sus 7 métricas}}{7}
  \qquad
  \text{Escalabilidad} =
  \frac{\text{suma de sus 9 métricas}}{9}
\]
\[
  \text{Soporte y mantenimiento} =
  \frac{\text{suma de sus 4 métricas}}{4}
  \qquad
  \text{Adopción y madurez} =
  \frac{\text{suma de sus 3 métricas}}{3}
\]
\[
  \text{Observabilidad} =
  \frac{\text{suma de sus 7 métricas}}{7}
\]

Por ejemplo, si una herramienta obtiene calificaciones de 3, 3, 2 y 2 en
las cuatro métricas de soporte y mantenimiento, el valor de ese indicador
es $(3+3+2+2) \div 4 = 2{,}5$.

Los dos indicadores restantes, adecuación al contexto GRID y riesgo
tecnológico, no se obtienen con este promedio, sino a partir de una
valoración del ajuste operativo de la herramienta en el GRID y de los
estados proyectados para el año 2026.

\subsubsection{Paso 2: obtener la puntuación DAR}

A diferencia de las métricas, los indicadores \textbf{no} pesan todos lo
mismo: cada uno tiene asignado un porcentaje según su importancia para el
proyecto. Por eso la puntuación DAR no es un promedio simple, sino un
\textbf{promedio ponderado}: el valor de cada indicador se multiplica por
su porcentaje y luego se suman todos los resultados.

\begin{align*}
  \text{DAR} \;=\;\; & 0{,}25 \times \text{Rendimiento técnico}
                 \;+\; 0{,}20 \times \text{Escalabilidad} \\[0.2em]
       &+\; 0{,}20 \times \text{Soporte y mantenimiento}
        \;+\; 0{,}10 \times \text{Adopción y madurez} \\[0.2em]
       &+\; 0{,}10 \times \text{Observabilidad}
        \;+\; 0{,}10 \times \text{Adecuación al contexto GRID} \\[0.2em]
       &+\; 0{,}05 \times \text{Riesgo tecnológico}
\end{align*}

Los factores $0{,}25$, $0{,}20$, $0{,}10$ y $0{,}05$ son los porcentajes
de la lista anterior escritos en forma decimal ($25\,\%$, $20\,\%$,
$10\,\%$ y $5\,\%$), y entre todos suman el $100\,\%$ de la evaluación.
Así, un buen resultado en rendimiento técnico influye cinco veces más en
la puntuación final que un buen resultado en riesgo tecnológico.

Como todos los indicadores están entre 1 y 3 y los porcentajes suman el
$100\,\%$, la puntuación DAR también queda entre 1 y 3: un valor cercano
a 3 indica una herramienta líder frente a los criterios priorizados por
el GRID, mientras que un valor cercano a 1 señala una alternativa muy
limitada para las necesidades del proyecto. Esta puntuación es la que se
reporta en la última columna de las plantillas y la que permite ordenar
las herramientas en el análisis DAR.

La Tabla~\ref{tab:plantilla-proxy} presenta la plantilla de análisis DAR empleada para evaluar
las herramientas de proxy-caché, con los criterios y sus respectivos
pesos.

\begin{table}[H]
\centering\footnotesize
\begin{tabular}{|>{\raggedright\arraybackslash\bfseries}p{2.4cm}|c|c|c|c|c|c|c|c|c|}
\hline
\rowcolor{gridhead} \thead[l]{Criterios/\\Herram.} & \thead{Frec.\\SMS} & \thead{Rend.\\(25\%)} & \thead{Esc.\\(20\%)} & \thead{Sop.\\(20\%)} & \thead{Adop.\\(10\%)} & \thead{Obs.\\(10\%)} & \thead{Adec.\\(10\%)} & \thead{Riesgo\\(5\%)} & \thead{DAR} \\ \hline
HAProxy & 0 &  &  &  &  &  &  &  &  \\ \hline
\rowcolor{gridrow} NGINX & 1-2 &  &  &  &  &  &  &  &  \\ \hline
Envoy & 1-2 &  &  &  &  &  &  &  &  \\ \hline
\rowcolor{gridrow} Varnish & 0 &  &  &  &  &  &  &  &  \\ \hline
Traefik & 0 &  &  &  &  &  &  &  &  \\ \hline
\rowcolor{gridrow} Caddy & 0 &  &  &  &  &  &  &  &  \\ \hline
Apache Traffic Server & 0 &  &  &  &  &  &  &  &  \\ \hline
\rowcolor{gridrow} Kong & 1 &  &  &  &  &  &  &  &  \\ \hline
Squid & 0 &  &  &  &  &  &  &  &  \\ \hline
\rowcolor{gridrow} Istio & 0 &  &  &  &  &  &  &  &  \\ \hline
Linkerd & 0 &  &  &  &  &  &  &  &  \\ \hline
\rowcolor{gridrow} K8Sidecar & 1 &  &  &  &  &  &  &  &  \\ \hline
Proxy-Service & 1 &  &  &  &  &  &  &  &  \\ \hline
\rowcolor{gridrow} nuster & 0 &  &  &  &  &  &  &  &  \\ \hline
\end{tabular}
\caption{Plantilla análisis DAR proxy-caché}
\label{tab:plantilla-proxy}
\end{table}

De igual manera, en la Tabla~\ref{tab:plantilla-almac} se muestra la plantilla de análisis DAR
aplicada a las herramientas de almacenamiento.

\begin{table}[H]
\centering\footnotesize
\begin{tabular}{|>{\raggedright\arraybackslash\bfseries}p{2.7cm}|c|c|c|c|c|c|c|c|c|}
\hline
\rowcolor{gridhead} \thead[l]{Herramienta} & \thead{Frec.\\SMS} & \thead{Rend.\\(25\%)} & \thead{Esc.\\(20\%)} & \thead{Sop.\\(20\%)} & \thead{Adop.\\(10\%)} & \thead{Obs.\\(10\%)} & \thead{Adec.\\(10\%)} & \thead{Riesgo\\(5\%)} & \thead{DAR} \\ \hline
Ceph & 0 &  &  &  &  &  &  &  &  \\ \hline
\rowcolor{gridrow} Lustre & 0 &  &  &  &  &  &  &  &  \\ \hline
JuiceFS & 0 &  &  &  &  &  &  &  &  \\ \hline
\rowcolor{gridrow} DAOS & 0 &  &  &  &  &  &  &  &  \\ \hline
SeaweedFS & 0 &  &  &  &  &  &  &  &  \\ \hline
\rowcolor{gridrow} ROOT / RDataFrame / Spark / Dask / OSCAR & 1 &  &  &  &  &  &  &  &  \\ \hline
HDFS (EC + ISA-L / SYCL) & 1 &  &  &  &  &  &  &  &  \\ \hline
\rowcolor{gridrow} Knative + x264 / AV1 & 1 &  &  &  &  &  &  &  &  \\ \hline
Longhorn & 0 &  &  &  &  &  &  &  &  \\ \hline
\rowcolor{gridrow} ExaHDF5 & 1 &  &  &  &  &  &  &  &  \\ \hline
HDF5 & 1-2 &  &  &  &  &  &  &  &  \\ \hline
\rowcolor{gridrow} MinIO & 0 &  &  &  &  &  &  &  &  \\ \hline
HDFS & 1 &  &  &  &  &  &  &  &  \\ \hline
\rowcolor{gridrow} RustFS & 0 &  &  &  &  &  &  &  &  \\ \hline
BeeOND & 1 &  &  &  &  &  &  &  &  \\ \hline
\rowcolor{gridrow} OpenStack Swift & 0 &  &  &  &  &  &  &  &  \\ \hline
Garage & 0 &  &  &  &  &  &  &  &  \\ \hline
\rowcolor{gridrow} IPFS & 1 &  &  &  &  &  &  &  &  \\ \hline
GekkoFS & 1 &  &  &  &  &  &  &  &  \\ \hline
\rowcolor{gridrow} GlusterFS & 0 &  &  &  &  &  &  &  &  \\ \hline
BurstFS & 1 &  &  &  &  &  &  &  &  \\ \hline
\rowcolor{gridrow} TEES & 1 &  &  &  &  &  &  &  &  \\ \hline
\end{tabular}
\caption{Plantilla análisis DAR almacenamiento}
\label{tab:plantilla-almac}
\end{table}

\input{secciones/04-analisis-dar}
%% ====================================================================
%%  Tecnología ganadora
%% ====================================================================
\section{Tecnología ganadora}\label{sec:ganadora}

\begin{enumerate}[label=\textbf{\arabic*.}, leftmargin=*, itemsep=0.9em]
  \item \textbf{Proxy-Caché / Gateway}

  Del análisis comparativo realizado mediante la metodología DAR, se establece una arquitectura híbrida para la capa de proxy-caché del grupo de investigación GRID, combinando las fortalezas de dos tecnologías líderes:

  HAProxy como Proxy Inverso y Gestión de Certificados. HAProxy se posiciona como la tecnología ganadora para la función de proxy inverso y terminación TLS dentro del GRID, al obtener la puntuación DAR más alta del ranking (2.87/3.0). Es muy eficiente en CPU y RAM ($<$1 ms de latencia, 42k RPS, consumo CPU $\sim$50 \%), respaldada por benchmarks oficiales y una LTS garantizada hasta 2031. Para el GRID, HAProxy cumple de manera nativa y excelente las NPO críticas de automatización de certificados SSL (Let's Encrypt integrado), gobernanza del tráfico mediante ACLs granulares, rate limiting avanzado y adaptación a cambios mediante runtime API dinámica sin downtime. Su rol como proxy inverso garantiza la publicación segura de servicios internos hacia la comunidad académica, eliminando advertencias de navegadores y asegurando la confianza en los sitios web del grupo.

  NGINX como Forward Proxy y Control de Tráfico. NGINX se selecciona como la alternativa documentada y complementaria para las funciones de forward proxy, control de tráfico y aceleración en la descarga de archivos grandes. Aunque su puntuación DAR global (2.49) es inferior a HAProxy, NGINX aporta una documentación académica extensa y una caché HTTP robusta validada por F5 en 2026, aspectos relevantes para el contexto universitario del GRID. En su rol de forward proxy, NGINX permite la aceleración de descargas de ISOs y contenido académico mediante caché HTTP con cache hit ratio superior al 80 \%, mientras que sus capacidades de QoS y ACLs facilitan el control del uso no misional de la red. Su integración con la infraestructura existente del GRID es plug-and-play, minimizando la curva de aprendizaje y sirviendo como respaldo documentado ante cualquier escenario de fallo o migración futura.

  \item \textbf{Almacenamiento de Archivos de Gran Tamaño}

  Para el backend de almacenamiento de objetos grandes (ISOs, VMs, datasets) del GRID, el análisis DAR identifica una estrategia escalonada que equilibra la producción con la agilidad del prototipo funcional requerido en el trabajo de grado:

  Ceph como Backend Principal de Producción: Ceph lidera el ranking de almacenamiento con una puntuación DAR de 2.89/3.0, siendo la tecnología ganadora indiscutible para el repositorio local robusto del GRID. Su arquitectura distribuida con replicación 3x + erasure coding (EC) y auto-healing mediante el algoritmo CRUSH, garantiza la persistencia y tolerancia a fallos de los archivos grandes. Con más de 500K IOPS en NVMe (BlueStore) y respaldo enterprise de IBM/Red Hat (versiones Reef, Squid y Tentacle activas en 2026), Ceph satisface la necesidad de un repositorio local de archivos grandes, escalabilidad sin reemplazo de hardware y alta disponibilidad. Su integración con Prometheus y Grafana permite el monitoreo en tiempo real del consumo de almacenamiento y salud del cluster, alineándose con los criterios de observabilidad del proyecto.

  SeaweedFS como Backend del Prototipo Funcional: SeaweedFS se selecciona como la tecnología ganadora para el prototipo funcional del trabajo, al ofrecer un equilibrio óptimo entre funcionalidad S3-compatible y menor complejidad operativa respecto a Ceph. Con una puntuación DAR de 2.49, SeaweedFS soporta erasure coding nativo, API S3 completa y tiempos de respuesta de 2.1 ms para archivos pequeños, siendo suficiente para validar el marco operacional de proxy-caché del GRID sin requerir la curva de aprendizaje y recursos de un cluster Ceph completo. Su soporte activo por parte de Chris Lu (v3.85+) y su arquitectura liviana permiten desplegar el prototipo en la infraestructura no productiva del grupo de investigación de manera ágil, cumpliendo con el alcance del anteproyecto y generando métricas reales de throughput y cache hit ratio.

  Ceph destaca en producción a gran escala, pero su arquitectura implica una complejidad operativa excesiva para el alcance de un prototipo funcional. SeaweedFS ofrece una API S3 completa, acceso a disco O(1) y una arquitectura ligera que se despliega en minutos sobre la infraestructura del GRID, con una huella de CPU y memoria muy inferior. Además, incorpora métricas Prometheus nativas y dashboards Grafana, junto con erasure coding y replicación configurables, aportando observabilidad y tolerancia a fallos suficientes para validar el marco operacional proxy-caché sin la curva de aprendizaje ni los requisitos de hardware de un clúster Ceph. Por ello, tras reponderar la matriz DAR con criterios centrados en el prototipo, SeaweedFS supera a Ceph tanto en la valoración individual (2,83 frente a 2,70) como en la media consolidada del equipo (2,66 frente a 2,63), por lo que se selecciona como backend del prototipo.

\end{enumerate}

%% ====================================================================
%%  Referencias
%% ====================================================================
\section{Referencias}\label{sec:referencias}

CMMI Institute. (2010). \emph{CMMI \textregistered{} para Desarrollo, Versión 1.3 Equipo
del Producto CMMI}. \emph{1}(1), 1--555. \url{http://www.sei.cmu.edu}


\end{document}
